\section*{第 \ref{ch:b1:elementary-steps} 章——反应机理：基元反应与近似}

\begin{solution}{exo:b1:elementary-steps:1}
(a) 单分子：可以是\omterm{def:b1:elementary-steps:elementary-step}{基元反应}。(b) 双分子：可以是\omterm{def:b1:elementary-steps:elementary-step}{基元反应}。
(c) 三个分子、四个键同时重组：不是\omterm{def:b1:elementary-steps:elementary-step}{基元反应}。
(d) 三分子：由 \ce{N2} 带走所释放的能量，是少见但确实存在的
\omterm{def:b1:elementary-steps:elementary-step}{基元反应}。
\end{solution}

\begin{solution}{exo:b1:elementary-steps:2}
$v = k[\mathrm{A}][\mathrm{B}]$; $v = k[\mathrm{A}]^2$; $v = k[\mathrm{A}]$.
总反应方程式只是一份“收支表”，并不说明分子
如何相遇；其\omterm{def:b1:rate-laws:order}{速率方程}必须通过实验测定，或由
\omterm{def:b1:elementary-steps:elementary-step}{反应机理}推导得出。
\end{solution}

\begin{solution}{exo:b1:elementary-steps:3}
\omterm{def:b1:elementary-steps:elementary-step}{反应中间体}对应 30 处的谷；两个\omterm{def:b1:elementary-steps:energy-profile}{过渡态}分别对应
70 和 55 处的峰，其中最高的是第一个\omterm{def:b1:elementary-steps:energy-profile}{过渡态}：
第一步为\omterm{def:b1:elementary-steps:rds}{决速步}。该步骤为吸热过程（从 0 升至 30）。
\end{solution}

\begin{solution}{exo:b1:elementary-steps:4}
同样地，因子为 $10^4$（见\cref{prop:b1:elementary-steps:catalysis}）；
平衡常数不变；达到平衡约快
$10^4$ 倍。
\end{solution}

\begin{solution}{exo:b1:elementary-steps:5}
$t_{\max} = \ln(0.30/0.10)/(0.30 - 0.10) = \qty{5.49}{s}$。则
$[\mathrm{B}]_{\max} = a_0\frac{0.10}{0.20}(\eu^{-0.549} - \eu^{-1.648}) =
0.5\,a_0(0.577 - 0.192) = 0.192\,a_0$.
\end{solution}

\begin{solution}{exo:b1:elementary-steps:6}
$\dd[\mathrm{I}]/\dd t = k_1[\mathrm{A}] - (k_{-1} + k_2)[\mathrm{I}] = 0$
由此得 $[\mathrm{I}] = k_1[\mathrm{A}]/(k_{-1} + k_2)$，且 $v = k_2[\mathrm{I}]
= \frac{k_1k_2}{k_{-1}+k_2}[\mathrm{A}]$。若 $k_2 \gg k_{-1}$，则 $v =
k_1[\mathrm{A}]$，即第一步为\omterm{def:b1:elementary-steps:rds}{决速步}。若 $k_2 \ll k_{-1}$，
则 $v = (k_1/k_{-1})k_2[\mathrm{A}]$，即先有前置平衡，随后
是一个慢步。
\end{solution}

\begin{solution}{exo:b1:elementary-steps:7}
总反应：\ce{2H2O2 -> 2H2O + O2}。\omterm{def:b1:elementary-steps:catalyst}{催化剂}：\ce{I-}（先被消耗，随后
再生）；\omterm{def:b1:elementary-steps:elementary-step}{反应中间体}为 \ce{IO-}。慢的第一步决定了
\omterm{def:b1:rate-laws:order}{速率方程}：$v = k_1[\ce{H2O2}][\ce{I-}]$。
\end{solution}

\begin{solution}{exo:b1:elementary-steps:8}
该步骤能量明显升高：根据哈蒙德假设，其\omterm{def:b1:elementary-steps:energy-profile}{过渡态}
与碳正离子相似。能使碳正离子能量降低的取代基，
同样几乎使\omterm{def:b1:elementary-steps:energy-profile}{过渡态}能量降低，因此
能垒下降，该步骤变快。
\end{solution}

\begin{solution}{exo:b1:elementary-steps:9}
$k_{\mathrm B}/k_{\mathrm C} = \exp(20000/RT)$：在 \qty{300}{K} 时为 \num{3.0e3}，
在 \qty{600}{K} 时为 $\exp(4.01) = 55$。升温削弱了动力学上对
B 的偏好，并通过使 B 的生成成为可逆过程，
使体系达到更稳定的 C。
\end{solution}

\begin{solution}{exo:b1:elementary-steps:10}
$k_1[\mathrm{A}][\mathrm{M}] = (k_{-1}[\mathrm{M}] + k_2)[\mathrm{A}^*]$,
于是 $v = k_2[\mathrm{A}^*] = \dfrac{k_1k_2[\mathrm{A}][\mathrm{M}]}{k_{-1}[\mathrm{M}]
+ k_2}$。高压下（$k_{-1}[\mathrm{M}] \gg k_2$），$v =
(k_1k_2/k_{-1})[\mathrm{A}]$，为一级；低压下，$v =
k_1[\mathrm{A}][\mathrm{M}]$，为二级；碰撞活化成为
慢步。
\end{solution}

\begin{solution}{exo:b1:elementary-steps:11}
证明与定理中的相同。对 $k_1 = k_2 = k$：
$\frac{\dd}{\dd t}([\mathrm{B}]\eu^{kt}) = ka_0$，于是 $[\mathrm{B}]\eu^{kt} =
ka_0t$，得 $[\mathrm{B}] = a_0kt\,\eu^{-kt}$，在 $t = 1/k$ 时取最大值。
\end{solution}

\begin{solution}{exo:b1:elementary-steps:12}
$\mathrm{A} + \mathrm{B} \rightleftharpoons \mathrm{I} + \mathrm{C}$（快，
平衡常数 $K_1$），随后 $\mathrm{I} + \mathrm{B} \to \mathrm{P}$（慢，
$k_2$）。二者相加即为总反应；$[\mathrm{I}] =
K_1[\mathrm{A}][\mathrm{B}]/[\mathrm{C}]$，且 $v = k_2[\mathrm{I}][\mathrm{B}]
= k_2K_1[\mathrm{A}][\mathrm{B}]^2/[\mathrm{C}]$：副产物 C 在
前置平衡中生成，会减慢反应。
\end{solution}

\begin{solution}{pb:b1:elementary-steps:1}
\textbf{1.} 加倍 $[\ce{NO}]_0$，$v_0$ 变为 4 倍：二级；
$[\ce{O2}]_0$ 加倍则速率变为 2 倍：一级。
\textbf{2.} \omterm{def:b1:rate-laws:order}{速率方程}为 $v = k[\ce{NO}]^2[\ce{O2}]$；$k = 1.0 \times 10^{-5}/((1.0
\times 10^{-3})^2 \times 1.0 \times 10^{-3}) =
\qty{1.0e4}{L^2.mol^{-2}.s^{-1}}$.
\textbf{3.} 三个分子同时碰撞非常少见，而
一个单步反应不可能随升温而变慢（见第 5 题）。
\textbf{4.} $E_a = R\ln(k_{400}/k_{300})/(1/300 - 1/400) = 8.314 \ln(0.1)
/(8.33 \times 10^{-4}) = \qty{-23}{kJ/mol}$.
\textbf{5.} \omterm{def:b1:elementary-steps:elementary-step}{基元反应}通过越过能垒的碰撞而进行；
能够越过能垒的碰撞比例为 $\eu^{-E_a/RT}$，其中 $E_a \ge
0$，只会随温度升高而增大。
\textbf{6.} $-\dd[\ce{NO}]/\dd t = 2v$, $-\dd[\ce{O2}]/\dd t = v$.
\textbf{7.} \ce{2NO -> N2O2} 与 \ce{N2O2 + O2 -> 2NO2} 相加，
得总反应 \ce{2NO + O2 -> 2NO2}；\omterm{def:b1:elementary-steps:elementary-step}{反应中间体}为 \ce{N2O2}。
\textbf{8.} 两者均为双分子（第一步的逆反应为单分子）。
\textbf{9.} $v = k_2[\ce{N2O2}][\ce{O2}]$.
\textbf{10.} $[\ce{N2O2}] = K_1[\ce{NO}]^2$，其中 $K_1 = k_1/k_{-1}$。
\textbf{11.} $v = k_2K_1[\ce{NO}]^2[\ce{O2}]$：即观测到的\omterm{def:b1:rate-laws:order}{速率方程}。
\textbf{12.} $k = k_1k_2/k_{-1}$.
\textbf{13.} $\dd[\ce{N2O2}]/\dd t = k_1[\ce{NO}]^2 - k_{-1}[\ce{N2O2}] -
k_2[\ce{N2O2}][\ce{O2}]$.
\textbf{14.} 令其为零，得 $[\ce{N2O2}] = k_1[\ce{NO}]^2/(k_{-1} +
k_2[\ce{O2}])$.
\textbf{15.} 由 $v = k_2[\ce{N2O2}][\ce{O2}]$ 即得所述\omterm{def:b1:rate-laws:order}{速率方程}。
\textbf{16.} $k_{-1} \gg k_2[\ce{O2}]$：二聚体分解的次数远多于
与氧气相遇的次数，因此第一步保持平衡。
\textbf{17.} 若 $k_2[\ce{O2}] \gg k_{-1}$，则 $v = k_1[\ce{NO}]^2$：
对 \ce{NO} 为二级，对 \ce{O2} 为零级。
\textbf{18.} 实测对 \ce{O2} 为一级，对应第一种极限。
\textbf{19.} $\ln k = \ln k_1 + \ln k_2 - \ln k_{-1}$.
\textbf{20.} 对 $T$ 求导，并利用 $\dd\ln k_i/\dd T
= E_i/RT^2$，得 $E_a = E_1 + E_2 - E_{-1}$。
\textbf{21.} 升温使慢步略微加快（$E_2$ 较小），但使
前置平衡强烈地向 \ce{NO} 一侧移动（$E_{-1}$ 较大）：二聚体减少，
总速率随之下降。
\textbf{22.} 二聚体的能量比反应物低 $E_1 - E_{-1} = \qty{-38}{kJ/mol}$，
第二\omterm{def:b1:elementary-steps:energy-profile}{过渡态}高于二聚体 $E_2 = \qty{15}{kJ/mol}$，
因此比反应物低 \qty{23}{kJ/mol}。路径的最高点是
第一\omterm{def:b1:elementary-steps:energy-profile}{过渡态}，仅高出 \qty{2}{kJ/mol}。
\textbf{23.} $\exp\big(\frac{23000}{8.314}(\frac{1}{400} - \frac{1}{300})\big)
= 0.10$：在 \qty{400}{K} 时慢了十倍。
\textbf{24.} $E_a = 2 - 40 + 15 = \textbf{\qty{-23}{kJ/mol}}$，
即第一部分中测得的数值：该机理既解释了\omterm{def:b1:rate-laws:order}{速率方程}，
也解释了它那奇特的温度依赖性。
\end{solution}
